\documentclass[11pt]{article}
\usepackage[a4paper,margin=1in]{geometry}
\usepackage{fontspec}
\usepackage[boldfont=false]{xeCJK}
\setCJKmainfont{FandolSong-Regular.otf}
\usepackage{amsmath}
\usepackage{amssymb}
\usepackage{array}
\usepackage{booktabs}
\usepackage{graphicx}
\usepackage{adjustbox}
\usepackage{enumitem}
\setlist[itemize]{topsep=2pt,partopsep=0pt,parsep=0pt,itemsep=2pt,leftmargin=1.4em}
\setlist[enumerate]{topsep=2pt,partopsep=0pt,parsep=0pt,itemsep=2pt,leftmargin=1.6em}
\usepackage{parskip}
\usepackage{fvextra}
\fvset{breaklines=true,breakanywhere=true,fontsize=\small}
\setlength{\parindent}{0pt}

\begin{document}

\begin{center}
  {\LARGE\bfseries Infectious diseases}\\[0.35em]
  {\large A Level Biology}
\end{center}
\vspace{0.6em}

\section*{What causes infectious disease}

An \textbf{infectious disease} 传染病 is caused by a \textbf{pathogen} 病原体 — an organism that lives in or on a host and causes harm. Infectious diseases are \textbf{transmissible}: the pathogen can be \textbf{transmitted} 传播 (passed) from one person to another.

You need to know four diseases, the pathogen that causes each, and how each spreads.

\begin{center}
\adjustbox{max width=\linewidth}{%
\begin{tabular}{lll}
\toprule
\textbf{Disease} & \textbf{Pathogen and type} & \textbf{How it spreads} \\
\midrule
\textbf{cholera} 霍乱 & the \textbf{bacterium} 细菌 \emph{Vibrio cholerae} & drinking water or food \textbf{contaminated} 污染 with \textbf{faeces} 粪便 (human waste) \\
\textbf{malaria} 疟疾 & the \textbf{protoctist} 原生生物 \emph{Plasmodium} & the bite of an infected \textbf{mosquito} 蚊子, which acts as a \textbf{vector} 媒介 (a carrier of the pathogen); also through infected blood \\
\textbf{tuberculosis (TB)} 结核病 & the bacterium \emph{Mycobacterium} & tiny airborne \textbf{droplets} 飞沫 from coughs and sneezes; spreads fast where people are crowded \\
\textbf{HIV/AIDS} & the \textbf{virus} 病毒 HIV (which leads to \textbf{AIDS} 艾滋病) & unprotected sex, infected blood (for example shared needles), and from mother to baby \\
\bottomrule
\end{tabular}}
\end{center}

HIV infects and destroys certain white blood cells, so it slowly weakens the body's \textbf{immune system} 免疫系统.

\section*{Preventing and controlling these diseases}

Control has \textbf{biological}, \textbf{social} and \textbf{economic} sides — the science of the pathogen, people's behaviour and education, and the money and resources available. Examples:

\begin{itemize}
\item \textbf{cholera}: provide clean water and proper sewage treatment; good hygiene; \textbf{vaccines} 疫苗 in some areas.

\item \textbf{malaria}: sleep under nets; remove pools of still water where mosquitoes breed; spray \textbf{insecticides} 杀虫剂; take anti-malarial drugs.

\item \textbf{TB}: find and treat infected people with a long course of antibiotics; give the BCG vaccine; reduce overcrowding; trace contacts of patients.

\item \textbf{HIV}: use condoms; use clean needles; test donated blood; educate people. There is no cure and no vaccine yet, but drugs can slow the virus down.

\end{itemize}

In every case, cost (economic), people's willingness to change behaviour (social) and the supply of drugs or vaccines (biological) all affect how well a disease can be controlled.

\section*{Antibiotics}

An \textbf{antibiotic} 抗生素 is a drug that kills bacteria or stops them growing. For example, \textbf{penicillin} 青霉素 stops bacteria from building their \textbf{cell walls} 细胞壁. As the bacterium grows, its weak wall cannot hold it, so the cell takes in water and bursts.

Antibiotics do \textbf{not} work against \textbf{viruses}. A virus has no cell wall and no chemical reactions of its own to attack — it simply uses the machinery of the host cell. So there is no antibiotic target in a virus.

\section*{Antibiotic resistance}

Sometimes a \textbf{mutation} 突变 makes a bacterium \textbf{resistant} to an antibiotic, which means the antibiotic no longer kills it. This \textbf{resistance} 耐药性 is a serious problem:

\begin{itemize}
\item when an antibiotic is used, the non-resistant bacteria die, but any resistant ones survive and multiply. Over time, more and more bacteria carry the resistance.

\item some infections then become very hard, or impossible, to treat.

\end{itemize}

Steps to slow resistance down:

\begin{itemize}
\item only use antibiotics when they are really needed (not for viral illnesses such as colds).

\item always finish the full course, so no bacteria are left alive.

\item use the correct antibiotic for the infection.

\item reduce the heavy use of antibiotics in farming.

\end{itemize}


\end{document}
